\subsection{Literature Synthesis and Research Gap}
\label{subsec:literature-synthesis-gap}

The reviewed literature supports a calibrated gap. State-of-the-art stenosis
modeling spans resolved 3D/FSI simulations, geometrical multiscale coupling,
distributed 1D reductions, open solver workflows, well-balanced and high-order
schemes, and emerging surrogate or physics-informed learning methods. These
families answer different questions. Resolved models expose local velocity,
pressure, and traction fields at high setup cost. Reduced models support
repeated studies only after their closure and observation maps are declared.
Well-balanced and high-order schemes improve numerical evidence only for the
equilibria and outputs they actually test. Learned surrogates reduce evaluation
cost only after training data, uncertainty, and admissible operating regimes are
made part of the evidence record.

The present work addresses a practical reproducibility and observation-operator
gap for this implementation of an extended 1D stenosis model. It specifies the
solver-coordinate map, the Canic-derived model variant, targeted verification
evidence, and a declared plane--tetrahedron velocity comparison operator. Its
claim is therefore intentionally narrower than a field-wide state-of-the-art
claim. A stronger novelty or validation claim would require broader evidence
that comparable auditable implementations and matched 1D--3D operators do not
already exist, together with matched pressure, experimental, or clinical data.
That broader claim is outside this report.
